% Section 4. Owners: Tejas (BOINC), Kushagra (SWIM, VECA), Abhishek (Bully, VECTrust).
\section{Literature Review}
\label{sec:literature}

\subsection{Volunteer Computing with BOINC}
\todo{Tejas, about 0.8 page, on \cite{anderson2004boinc} and \cite{anderson2020boinc}: the project server model, work units and replication, result validation by comparing replicas (and its homogeneous redundancy for floating-point results), credit, and the central server as a single point of failure. End with what SpareCycles takes from BOINC and what it changes.}

\subsection{Group Membership and Failure Detection: SWIM}
Large distributed systems need every member to learn, quickly and at low cost, that another member has failed. The classic approach, in which every member periodically sends a heartbeat to every other member, produces a network load that grows quadratically with the group size, and either detection time or false-positive rate degrades as the group grows. Das, Gupta and Motivala~\cite{das2002swim} address this with SWIM, whose central idea is to separate \emph{failure detection} from the \emph{dissemination} of membership updates.

For detection, each member runs a protocol period of fixed length. In every period it picks one other member and sends it a \textit{ping}. If no acknowledgement arrives within a timeout, it does not declare the member failed; instead, it asks $k$ other members to probe the target on its behalf with a \textit{ping-req}. Only if none of these indirect probes succeeds is the target considered failed. The indirect probes make the detector robust to a single congested or lossy link between two healthy members, and the authors show that the probability of a false positive falls quickly as $k$ grows. Because each member sends a constant number of messages per period, the load per member is independent of the group size, and the expected time until some member first detects a failure is a constant number of protocol periods, $e/(e-1)$, also independent of the group size.

For dissemination, SWIM does not broadcast. Membership updates are \emph{piggybacked} on the ping, ping-req and acknowledgement messages that the detector already sends, so they spread in the manner of an epidemic, reaching the whole group in a number of periods that grows only logarithmically with its size, at no extra message cost.

The paper adds two refinements that our implementation adopts. First, a \emph{suspicion} subprotocol: a member that does not answer is first marked \textit{suspect} rather than dead, and the suspicion is gossiped to the group. If the suspected member is alive, it learns of the suspicion and refutes it by announcing itself alive with a higher \emph{incarnation number}; only a suspicion that survives a timeout becomes a confirmed failure. Incarnation numbers give a precedence order between conflicting updates about the same member, which keeps the protocol consistent without coordination. Second, probe targets are chosen in a randomly shuffled \emph{round-robin} order rather than uniformly at random, which bounds the worst-case time until a failed member is probed by some member to a number of periods linear in the group size.

SWIM is a natural fit for volunteer computing. Volunteers join and leave continuously, the group can grow large, and a coordinator that relies on lease expiry alone learns of a crash only after a full lease, which for our 60-second units is three minutes. With SWIM, the coordinator learns of the crash within seconds and reassigns the lost work early (Section~\ref{sec:scenario-volunteer}).

\subsection{Volunteer Edge-Cloud Computing: VECA}
Yeddulapalli et~al.~\cite{yeddulapalli2024veca} study volunteer edge-cloud computing, in which individuals and organizations contribute laptops, desktops and servers to run machine-learning and deep-learning workflows. They identify two obstacles: the nodes are intermittent and volatile, which makes workflow performance unpredictable, and the owners of the nodes must not see the workflow's models and data. Their framework, VECA, makes a three-fold contribution.

The first is \emph{capacity-based clustering}: nodes are grouped with $k$-means on their number of CPUs, memory and storage, and the clustering is repeated when the number of nodes grows by 10\%, so that a scheduler searches one suitable cluster instead of all nodes. The second is a \emph{two-phase scheduler}. A cloud hub first selects the cluster whose capacity best matches the workflow; an agent in that cluster then uses a recurrent neural network to forecast each node's probability of being online, ranks the nodes, and schedules on nodes whose predicted availability is at least 0.8, preferring the one nearest the user. If the chosen node fails, a fail-over step reads the stored workflow details and the precomputed node order from a Redis cache and resumes execution on the next node, without running the forecasting model again. The third is \emph{confidential computing}: confidential workflows run only on nodes that support AWS Nitro Enclaves, using an encrypted image, isolated CPU and memory, and an attestation document that proves the enclave's identity.

VECA was implemented with OpenFaaS, MicroK8s and Docker and evaluated in simulation with 50 nodes in four clusters, synthetic availability traces and two scientific workflows. It reports a two-fold reduction in node search latency over the next best approach and a mean productivity rate (the share of execution time not spent on recovery) of 86.9\%, against 66.7\% and 65.7\% for two earlier frameworks.

VECA and SpareCycles address the same root problem, volatile volunteer nodes, from opposite directions. VECA tries to \emph{avoid} failures by predicting availability and placing work on nodes that are likely to stay online, and it recovers by fail-over to the next node. SpareCycles assumes that failures and even malicious nodes are normal: it \emph{detects} failures by gossip, \emph{recovers} through leases and reassignment, and \emph{tolerates} incorrect results through replication and voting, in the BOINC tradition. VECA's availability forecasting would complement our resource discovery (Section~\ref{sec:design}), and its use of enclaves addresses a confidentiality requirement that SpareCycles deliberately leaves out of scope: our workloads are public scientific computations, and we protect the \emph{integrity} of tasks with signatures rather than their confidentiality.

\subsection{Leader Election: the Bully Algorithm}
\todo{Abhishek, about 0.4 page, on \cite{garciamolina1982elections}: assumptions (synchronous system, failure detection by timeout), the election, answer and coordinator messages, the highest live identifier wins, message complexity (worst case quadratic), and why it is unsafe under network partitions. Explain how term numbers and a majority leader lease fix this in SpareCycles.}

\subsection{Trusted Resource Allocation: VECTrust}
\todo{Abhishek, about 0.4 page, on \cite{pandey2021vectrust}: how VECTrust computes trust for volunteer resources and uses it in allocation, its evaluation, and how SpareCycles' resource discovery (trust scores updated from voting) builds on it.}
